%% kapa: exam=1 round=32 year=2021 period=2 form=A subject=appraisal_law subject_ko=감정평가관계법규 seq=34 number=34
%% source: Q-Net 감정평가사 기출문제 (cst00309) artlSeq=5208397 file=extracted/1008/5208397/A형/2021년도 제32회 감정평가사 1차 2교시 A형.pdf
%% conversion: kordoc 4.24.51 -> tools/gen1.py (auto); review=proofread-sonnet; flags=none
\begin{question}{34}
부동산등기법령상 신탁에 관한 등기에 관한 설명으로 옳은 것을 모두 고른 것은?
\begin{framedtext}
ㄱ. 수탁자가 여러 명인 경우 등기관은 신탁재산이 합유인 뜻을 기록하여야 한다. \par
ㄴ. 위탁자가 수탁자를 대위하여 신탁등기를 신청하는 경우 신탁등기의 신청은 해당 부동산에 관한 권리의 설정등기의 신청과 동시에 하여야 한다. \par
ㄷ. 수익자나 위탁자는 수탁자를 대위하여 신탁등기의 말소등기를 신청할 수 없다. \par
ㄹ. 법원은 수탁자 해임의 재판을 한 경우 지체 없이 신탁원부 기록의 변경등기를 등기소에 촉탁하여야 한다. \par
\end{framedtext}
\begin{choices}
  \item ㄱ, ㄷ
  \item ㄱ, ㄹ
  \item ㄴ, ㄷ
  \item ㄱ, ㄴ, ㄹ
  \item ㄴ, ㄷ, ㄹ
\end{choices}
\end{question}
\answer{2}
